\slajdid{02-s017}
\begin{frame}[label=02-s017]
\gusttekst
\[
F=(C+B+A)(C+B+\bar A)(C+\bar B+A)(\bar C+B+A)(\bar C+\bar B+\bar A)
\]
U oba slučaja dobijamo isti oblik funkcije koji je \alert{proizvod potpunih zbirova} i naziva se \alert{normalnom konjuktivnom formom} ili \alert{kanoničkom konjuktivnom formom}.

Gledajući oba načina, na koji smo došli do funkcije, lako je uočiti formalizam da se napiše ovakav oblik funkcije: Vrsta u kojoj funkcija ima vrednost logičke nule pojavljuje se kao zbir promenljivih formiranih tako da ako promenjiva ima vrednost logičke nule u zbiru se pojavljuje sa pravom vrednošću, a ako ima vrednost logičke jedinice u zbiru se pojavljuje sa komplementnom vrednošću. Ovako formirani zbirovi se vezuju proizvodima u formiranju izlazne funkcije.
\end{frame}
